\section{Discussion}

The results support the main implementation hypothesis only in the aggregate and with important qualifications. Adaptive beacon covariance improves beacon-only mean position RMSE relative to fixed covariance, and it improves the beacon-only velocity RMSE in all three scenarios. However, the adaptive beacon filter is slightly worse than the fixed-covariance beacon filter in Scenario 1 position RMSE. The evidence therefore favors adaptive covariance as a useful model for these saved data, not as a uniformly dominant choice for every trajectory and beacon geometry.

The fused LiDAR plus adaptive beacon case gives the lowest aggregate position RMSE, but the margin over LiDAR-only filtering is modest because the synthetic LiDAR pose measurements are frequent and usually available. This behavior is consistent with the availability table: beacon measurements are intermittent and often underdetermined, while LiDAR pose availability remains above 94\% in all scenarios.

The implementation has several limitations that affect interpretation. First, the filter initializes from the first truth state, so the experiments evaluate propagation and measurement-update behavior rather than global initialization. Second, the LiDAR backend uses synthetic pose measurements generated from truth and scan quality; raw scan matching is not implemented. Third, the RSSI model is a linear synthetic attenuation model, and the filter only adapts covariance from measured range and RSSI. Fourth, acceleration and gyro biases are injected into prediction inputs but are not included in the filter state. These choices are consistent with the repository scope, but they limit claims about real hardware performance.
